\section{Introduction}

The transition to post-quantum key agreement is underway. NIST has
standardized ML-KEM, other state-level bodies are providing their own
cryptographic recommendations, and organizations are beginning to update to
to protect against future quantum adversaries. However, these algorithms are
relatively new, and there remains a possibility of unexpected cryptanalytic
breakthroughs or vulnerable implementations~\cite{EC:CasDec23}. A natural
hedge against this risk is to deploy hybrid constructions that combine
post-quantum and classical algorithms, maintaining security as long as at
least one component remains secure.

This ``defense in depth'' approach is appealing: if the post-quantum KEM is
broken by classical cryptanalysis or a software bug, the classical component
provides classical security; if the classical component is broken by a
quantum computer, the post-quantum component provides quantum
resistance. However, constructing provably secure hybrid KEMs requires
care. The combiner function must preserve security properties under the right
assumptions, and different designs involve different trade-offs between the
cryptographic assumptions of the components and efficiency of the whole
scheme.

The Internet Research Task Force (IRTF)'s Crypto Forum Research Group (CFRG)
is working towards finalizing a recommendation for specific constructions of
hybrid KEMs. The current CFRG document~\cite{irtf-cfrg-hybrid-kems-07} gives
four constructions named CG, CK, UG, and UK. The CG and UK constructions are,
roughly speaking, covered by existing analyses of similar schemes. These
analyses do not apply to CK and UG. The goal of this paper is to close this
gap and provide a clean analysis of UG; a subset of this paper's authors are
also working in parallel on a similar paper for CK.

\paragraph{Our Contributions.} Our first contribution is to formally specify
the UG construction (Section~\ref{sec:ug}). At a high level, UG (hereafter
$\UG$) combines a post-quantum KEM with a classically-secure nominal
group. Like other hybrids, $\UG$ concatenates shared secrets, ciphertexts,
and encapsulation keys from both components, then applies a key derivation
function (KDF) to produce the final shared secret. Its design is overall
conservative---$\UG$ includes all public values in the KDF input. The main
departure from other hybrid constructions is that the group elements derived
during the DH key computation are put into the KDF directly instead of being
hashed once and then having the result hashed again in the final KDF. This
breaks an abstraction boundary and renders prior security analyses
inapplicable, but saves redundant hashing and eases implementation with
existing libraries.

The next contribution (Section~\ref{sec:ug-proofs}) is a clean, modular
analysis of the IND-CCA security of $\UG$. Like other hybrid KEMs, we prove
two complementary results corresponding to two different possible worlds. In
Section~\ref{sec:ug-proofs:sdh} we show that even if the KEM component is
broken, $\UG$ remains secure (modelling the KDF as a random oracle) as long
as the strong Diffie-Hellman problem in the nominal group holds. Our proof
roughly follows the analogous result from X-Wing~\cite{rfcdraft}, though
because $\UG$ hashes the KEM ciphertext we need not reduce to ciphertext
preimage resistance (C2PRI)~\cite{CiC:BCDKSV24}.

In Section~\ref{sec:ug-proofs:cca} we show that IND-CCA security of $\UG$
also holds when reducing \emph{only} to the IND-CCA of the component KEM and
the PRF security of the KDF. This proof largely follows the template for
proving CCA of KEM combiners given in~\cite{PKC:GiaHeuPoe18}. Notably, this
proof does not require modelling the KDF as a random oracle.  We end the
security analyses with a brief discussion of other security goals, like
key-binding~\cite{CCS:CreDaxMed24}.


We conclude the paper in Section~\ref{sec:concrete} by describing some
concrete choices of nominal group and KEM that can be used to instantiate
$\UG$. Since $\UG$ requires fairly minimal assumptions from its components,
it can be instantiated with diverse combinations of elliptic curves and KEMs.


%Notably, $\UG$ does not require ciphertext second-preimage resistance
%(C2PRI)\cite{CiC:BCDKSV24} from its components— SDH and IND-CCA security
%suffices, nor that the classical component be an IND-CCA KEM. Thus, $\UG$
%can be instantiated with diverse combinations of elliptic curves as nominal
%groups, and any IND-CCA secure post-quantum KEM.

Overall, our work complements and supports ongoing work on standardizing
hybrid KEM designs that strike desirable tradeoffs between performance and
security.  $\UG$ represents a conservative point in the design space: it
achieves strong security guarantees with minimal assumptions, at the cost of
hashing additional public data, while eliminating layers of abstraction.

%\paragraph{Structure of this paper.}
%Notation and definitions are introduced \autoref{sec:prelims}.  The $\UG$
%framework is introduced in \autoref{sec:ug}; our generic construction --
%named Heavy Quantum Bomber -- is a general framework that can be %instantiated
%with different component algorithms.  In \autoref{sec:ug-proofs} we prove
%the IND-CCA security when the different components fail in different ways.
